\subsubsection{Hurwitz--Siegel--Hecke 结构：\texorpdfstring{$S_3$}{S3} Hurwitz 叠、Klingen 型 \texorpdfstring{$3$}{3}-水平与 \texorpdfstring{$(3,3)$}{(3,3)}-同源对应}\label{par:xi-hurwitz-siegel-hecke-klingen-level3}
设 \(\mathcal H_{S_3}^{(2^6)}\) 为分支型别 \(((2,2,2,2,2,2))\) 的 \(S_3\)-Galois 覆盖 \(Y\to\PP^1\) 的连通 Hurwitz 叠，
其中分支点不标号且允许 \(\PP^1\) 的自同构。
记 \(\mathcal A_2\) 为主极化阿贝尔曲面之模叠，\(\mathcal J_2\subset \mathcal A_2\) 为 Jacobian 轨道（Torelli 嵌入之像）。
定义 Klingen 型 \(3\)-水平结构模叠
\[
\mathcal A_{2}^{\mathrm{Kl}}(3)
:=
\bigl\{(A,\lambda,\ell)\ \bigm|\ (A,\lambda)\in\mathcal A_2,\ \ell\subset A[3]\ \text{为一条 \(\FF_3\)-直线}\bigr\},
\]
其中 \(\ell\) 等价于一个阶 \(3\) 的循环子群。

\begin{theorem}[Hurwitz 模空间的 Klingen 实现与中间二次商的 Jacobian 解释]\label{thm:xi-hurwitz-klingen-jacobian-birational}
存在一个在稠密开子叠上为同构的双有理对应
\[
\mathcal H_{S_3}^{(2^6)}\ \dashrightarrow\ \mathcal A_{2}^{\mathrm{Kl}}(3)\big|_{\mathcal J_2}.
\]
其在几何点层面由如下构造给出：对 \(S_3\)-Galois 覆盖 \(Y\to\PP^1\)，令
\[
C:=Y/C_3,
\]
其中 \(C_3\triangleleft S_3\) 为唯一阶 \(3\) 正规子群。
则 \(C\to\PP^1\) 为分支点数为 \(6\) 的二次覆盖，从而 \(C\) 为属 \(2\) 的超椭圆曲线；
并且三次中间覆盖 \(Y\to C\) 为无分歧循环覆盖，等价于一个非零 \(3\)-挠类
\(
\alpha\in \Pic^0(C)[3]\setminus\{0\}
\)。
令 \(\ell:=\FF_3\cdot\alpha\subset J(C)[3]\)，则
\[
(Y\to\PP^1)\ \longmapsto\ \bigl(J(C),\lambda_C,\ell\bigr)
\]
给出所述对应，其中 \(\lambda_C\) 为 Jacobian 的主极化。
\end{theorem}
\begin{proof}[证明要点]
由 \(Y\to\PP^1\) 取 \(C_3\)-商得到 \(C\to\PP^1\) 的二次覆盖。
分支型别为六个换位元意味着每个分支点的惰性群阶为 \(2\)，故 \(C\to\PP^1\) 在六点处分歧；
Riemann--Hurwitz 公式给出 \(g(C)=2\)，从而 \(J(C)\) 定义一几何点落在 \(\mathcal J_2\) 上。
\par
另一方面，\(Y\to C\) 为无分歧循环三次覆盖，按循环无分歧覆盖与 \(\Pic^0(C)[3]\) 的对应，可由一非零 \(3\)-挠类 \(\alpha\) 唯一确定；
其张成的 \(\FF_3\)-直线 \(\ell=\FF_3\cdot \alpha\subset J(C)[3]\) 便给出 Klingen 型水平结构。
\par
反向方向在稠密开集上成立：当 \(C\) 的自同构群为 \(\{\pm1\}\) 且 \(J(C)\) 简单时，给定 \((J(C),\ell)\) 可复原 \(\alpha\)（在 \(\ell\setminus\{0\}\) 中选取任一生成元仅改变覆盖的同构），从而重建 \(Y\to C\)，再与超椭圆映射 \(C\to\PP^1\) 复合即可恢复 \(Y\to\PP^1\)。
因此在稠密开子叠上得到互为逆的构造，从而对应为双有理。证毕。
\end{proof}

\subsubsection{正交邻接与 \texorpdfstring{$(3,3)$}{(3,3)}-同源几何}
对 \((A,\lambda)\) 上的 Weil 配对在 \(A[3]\) 上诱导的辛型式，仍记为 \(\langle\cdot,\cdot\rangle_3\)。
对 \((A,\lambda,\ell)\in\mathcal A_2^{\mathrm{Kl}}(3)\) 定义正交补 \(\ell^\perp\subset A[3]\)，并设
\[
\mathcal T
:=
\bigl\{(A,\lambda,\ell,\ell')\ \bigm|\ (A,\lambda,\ell)\in\mathcal A_2^{\mathrm{Kl}}(3),\ \ell'\subset \ell^\perp,\ \ell'\neq \ell\bigr\}.
\]

\begin{proposition}[正交邻接对应与 \texorpdfstring{$(3,3)$}{(3,3)}-同源核]\label{prop:xi-orthogonal-adjacency-33-isogeny}
\leanverified{paper\_xi\_orthogonal\_adjacency\_33\_isogeny}
在稠密开集上，投影 \(\mathcal T\to\mathcal A_2^{\mathrm{Kl}}(3)\) 为有限 \'etale 且次数为 \(12\)。
并且对任意 \((A,\lambda,\ell,\ell')\in\mathcal T\)，令 \(\Delta:=\ell\oplus \ell'\subset A[3]\)，则 \(\Delta\) 为极大各向同性平面（Lagrangian \(\FF_3^2\)），从而诱导一条与主极化相容的 \((3,3)\)-同源
\[
A\ \longrightarrow\ A/\Delta.
\]
\end{proposition}
\begin{proof}
在 \(\FF_3\)-辛向量空间 \(V:=A[3]\cong \FF_3^4\) 中，\(\ell^\perp\) 为三维子空间。
其一维子空间总数为 \((3^3-1)/(3-1)=13\)，去除 \(\ell\) 本身得到 \(12\) 个可能的 \(\ell'\)，从而给出纤维次数 \(12\)；
在泛点处 \(A\) 简单且 \(A[3]\) 在模 \(3\) 单值化下无额外自同构时，对应为 \'etale。
\par
若 \(\ell'\subset\ell^\perp\) 且 \(\ell'\neq \ell\)，则 \(\langle x,y\rangle_3=0\) 对任意 \(x\in\ell,y\in\ell'\) 成立，故 \(\Delta=\ell\oplus\ell'\) 为二维各向同性子空间；
在维数 \(4\) 的辛空间中二维各向同性必为极大各向同性。
极大各向同性 \(\Delta\subset A[3]\) 与 \((3,3)\)-同源的核数据等价，从而得到所述同源（参见 \cite{BruinFlynnTesta2014Descent33Isogeny}）。证毕。
\end{proof}

\subsubsection{邻接算子谱与规范投影子}
记 \(f:\mathcal A_2^{\mathrm{Kl}}(3)\to\mathcal A_2\) 为忘却直线的有限态射，并记 \(\QQ^{40}:=f_\ast\QQ\) 为秩 \(40\) 的置换局部系统。
由命题 \ref{prop:xi-orthogonal-adjacency-33-isogeny} 的对应 \(\mathcal T\) 得到在 \(\QQ^{40}\) 上的自同态 \(T\)（邻接/Hecke 算子）。

\leanverified{paper\_xi\_klingen\_hecke\_charpoly\_projectors}
\begin{theorem}[Klingen 邻接算子的特征多项式与三重规范分解]\label{thm:xi-klingen-hecke-charpoly-projectors}
在任意几何纤维上，\(T\) 的最小多项式为
\[
(T-12)(T-2)(T+4)=0,
\]
其特征多项式为
\[
\chi_T(X)=(X-12)(X-2)^{24}(X+4)^{15}.
\]
因此可定义三个互补的规范幂等元
\[
e_{12}=\frac{(T-2)(T+4)}{160},\qquad
e_{2}=\frac{(12-T)(T+4)}{60},\qquad
e_{-4}=\frac{(T-12)(T-2)}{96},
\]
满足 \(e_{12}+e_{2}+e_{-4}=1\)，并分别投影到维数 \(1,24,15\) 的本征子空间。
\end{theorem}
\begin{proof}[证明要点]
几何纤维可识别为 \(\PP(A[3])\cong \PP^3(\FF_3)\) 的 \(40\) 个点（即 \(A[3]\) 的 \(\FF_3\)-直线）。
对应 \(\mathcal T\) 在纤维上给出“正交邻接图”：顶点为直线 \(\ell\)，邻点为 \(\ell^\perp\) 内不同于 \(\ell\) 的直线。
线性代数计数给出每点度数为 \(12\)。
该图是参数 \((v,k,\lambda,\mu)=(40,12,2,4)\) 的强正则图；
其谱性质（本征值 \(12,2,-4\) 与重数 \(1,24,15\)）由强正则图的一般理论确定，并可参见 \cite{BrouwerVanMaldeghemStronglyRegularGraphs}。
本征值互异时的 Lagrange 插值给出相应幂等投影公式。证毕。
\end{proof}

\subsubsection{旗标模叠与次数分解 \texorpdfstring{$12=4\cdot 3$}{12=4·3}}
记 \(\mathcal A_{2}^{\mathrm{Si}}(3)\) 为 Siegel 型 \(3\)-水平结构模叠，其几何点可表述为 \((A,\lambda,\Delta)\)，其中 \(\Delta\subset A[3]\) 为极大各向同性平面。
设
\[
\mathcal F
:=
\bigl\{(A,\lambda,\ell,\Delta)\ \bigm|\ (A,\lambda)\in\mathcal A_2,\ \ell\subset\Delta\subset A[3],\ \ell\ \text{直线},\ \Delta\ \text{Lagrangian}\bigr\},
\]
并令 \(\pi_{\mathrm{Kl}}:\mathcal F\to\mathcal A_{2}^{\mathrm{Kl}}(3)\) 忘却 \(\Delta\)，\(\pi_{\mathrm{Si}}:\mathcal F\to\mathcal A_{2}^{\mathrm{Si}}(3)\) 忘却 \(\ell\)。

\leanverified{paper\_xi\_flag\_factorization\_12\_4x3}
\begin{proposition}[正交邻接对应的旗标分解]\label{prop:xi-flag-factorization-12=4x3}
在稠密开集上，\(\pi_{\mathrm{Kl}}\) 为有限 \'etale 且次数为 \(4\)。
并且对固定的 \((A,\lambda,\Delta)\)，平面 \(\Delta\) 内除去给定直线 \(\ell\) 之外尚有 \(3\) 条直线；
从而由 \(\mathcal F\) 诱导的复合对应
\[
\mathcal A_{2}^{\mathrm{Kl}}(3)\ \xleftarrow{\ \pi_{\mathrm{Kl}}\ }\ \mathcal F\ \xrightarrow{\ \mathrm{swap}\ }\ \mathcal F\ \xrightarrow{\ \pi_{\mathrm{Kl}}\ }\ \mathcal A_{2}^{\mathrm{Kl}}(3)
\]
在几何纤维上与 \(\mathcal T\) 的正交邻接一致，并给出次数分解 \(12=4\cdot 3\)。
\end{proposition}
\begin{proof}
在 \(V\simeq\FF_3^4\) 的辛几何中，固定直线 \(\ell\) 后，商空间 \(\ell^\perp/\ell\) 为二维辛空间。
包含 \(\ell\) 的 Lagrangian 平面与 \(\ell^\perp/\ell\) 的一维子空间一一对应，故其数目为 \((3^2-1)/(3-1)=4\)，得到 \(\deg(\pi_{\mathrm{Kl}})=4\)。
而在固定二维平面 \(\Delta\simeq\FF_3^2\) 内，一维子空间总数同为 \(4\)，除去 \(\ell\) 余 \(3\) 条。
将两步复合即得 \(12\) 个邻接选择，且由构造可见对应与 \(\mathcal T\) 相同。证毕。
\end{proof}

\subsubsection{四十度覆盖的原始性与射影 \texorpdfstring{$3$}{3}-水平的 Galois 闭包}
设 \(\pi:\mathcal A_{2}^{\mathrm{Kl}}(3)\to\mathcal A_2\) 为忘却 \(\ell\) 的次数 \(40\) 有限态射，并将其限制到 \(\mathcal J_2\) 的稠密开集。

\begin{theorem}[原始性与 \texorpdfstring{$PSp(4,\FF_3)$}{PSp(4,F3)} 射影单群闭包]\label{thm:xi-40cover-primitive-galois-closure-psp43}
\(\pi\) 在上述开集上为原始覆盖：不存在非平凡因子分解 \(\pi=\pi_2\circ\pi_1\)（其中 \(\deg(\pi_i)>1\) 且中间空间正则）。
并且其 Galois 闭包在函数域层面同构于射影化扩张 \(\PP(A[3])\to \mathcal A_2\) 所给出的正则扩张；
其几何单群为 \(PSp(4,\FF_3)\)，对应于该群对 \(40\) 条 \(\FF_3\)-直线的自然射影作用。
\end{theorem}
\begin{proof}[证明要点]
在几何纤维上，\(\pi\) 等价于 \(PSp(4,\FF_3)\) 对 \(\PP^3(\FF_3)\) 的自然作用；
稳定子为“稳定一条各向同性直线”的抛物子子群，是极大子群（参见 \cite{Conrad2016StandardParabolicSubgroups}）。
置换作用的原始性与稳定子的极大性等价，从而得到覆盖的原始性。
\par
Galois 闭包对应于该作用的 core；对射影辛群在 \(\PP^3(\FF_3)\) 上的自然作用其 core 为平凡，
因此闭包群即为 \(PSp(4,\FF_3)\)。
该群与 \(A_2(3)^\ast\)（Igusa 紧化）中由 \(3\)-水平结构诱导的组合配置之间的对应，可参见 \cite{HoffmanWeintraubSp4F3AbelianSurfaces}。证毕。
\end{proof}

\endinput
